Marine autonomy requires robust perception across diverse conditions, often utilizing paired Infrared (IR) and RGB sensors. However, physical displacement and differing fields of view necessitate precise registration, while varying environmental conditions (fog, night, glare) demand adaptive data augmentation. Standard pipelines use fixed registration parameters and static augmentation schemes, which yield suboptimal results when conditions fluctuate—for example, rigid registration may fail during high-motion maneuvers, and standard color jitter is meaningless for thermal IR.

This proposal investigates a content-aware RL policy that jointly optimizes registration complexity and augmentation intensity. Unlike "search-then-freeze" methods like AutoAugment \cite{cubuk2019autoaugment}, our agent adapts *per-frame* or *per-sequence*, selecting simpler transforms for stable scenes and aggressive augmentation for low-diversity data. We hypothesize that this dynamic control will significantly outperform static baselines in varying marine environments.

Our contributions are: (1) a joint optimization framework for registration and augmentation; (2) an instance-level adaptation mechanism based on frame quality; and (3) a specific evaluation on multimodal marine datasets (M3FD, FLIR) targeting thermal crossover and motion artifacts. Fig.~
ef{fig:method-overview} illustrates the closed-loop pipeline, while Fig.~
ef{fig:rw-matrix} contrasts our approach with existing static and randomized methods.
